\section{Blessed are the Meek}

\begin{quotex}
Blessed are the meek: for they shall inherit the earth. (Matt. 5:5)
\end{quotex}

Those lacking in spiritual insight can only understand this beatitude in human terms. They understand meekness as something weak, timid, and passive. In contrast, they praise the earthy man with strong appetite and passionate nature. Both of these points of view reflect an inadequate notion of man. One reflects a Marxist view of an underclass motivated by resentment; the other, a modern neopagan caricature. When the real constitution of man is understood, it takes on a different meaning.

\paragraph{Powers of the soul}


Plato describes the three aspects or powers of the soul:

\begin{itemize}


\item \textbf{Intelligent aspect} (\emph{logikos}): this pertains to the intellect, which is the highest faculty in man by which he knows God or the essence of things by direct apprehension

\item \textbf{Appetitive aspect} (\emph{eros}): the soul's desiring power


\item \textbf{Incensive aspect} (\emph{thymos}): the force provoking vehement feelings, often anger. Another word for someone dominated by this aspect is ``spirited'', not to be confused with ``Spirit''
\end{itemize}


In the Chariot Allegory\footnote{See next paragraph.}, the soul is portrayed as a charioteer (the Intellect) controlling two horses. The white one represents the \emph{thymos} and the black one, the appetitive aspect. The latter is the aspect most susceptible to the influence of the passions and hence is more difficult to control. Even the former can erupt into useless anger.

Nevertheless, both lower aspects, including spiritedness, must be under the control of the intellect in order for the soul to continue its ascent. \textbf{Vladimir Solovyov}, in his \emph{Lectures on Divine Humanity} puts it this way: ``One can be considered free from passions only when one has them but has power over them, when one possesses, but is not possessed by them.''

In Matthew, the word translated as ``meek'' is \emph{praos}, which was also used to mean the gentling of a horse. The horse is the most noble of animals, but it must be broken to be of use to the knight\footnote{\url{https://gornahoor.net/?p=281}}.

So we see that the meek are those whose souls are under the control and guidance of the Intellect, since the appetitive and incensive aspects of the soul have been broken and put in service to the highest aspect in man. They still retain their driving forces: the appetitive to achieve a goal, and the incensive to provide the emotional energy and motivation to persist. Far from being mousy or timid, the meek are strong willed and inner directed. This is why they shall inherit the earth.

\paragraph{Parable of the Coach} This image represents the characteristics of man by a coach. The physical body is represented by the coach itself; the horses represent sensations, feelings and passions; the coachman is the ensemble of the intellectual faculties including reason; the person sitting in the coach is the master. In its normal state, the whole system is in a perfect state of operation: the coachman holds the reins firmly in his hands and drives the horses in the direction indicated by the master.

This, however, is not how things happen in the immense majority of cases. First of all, the master is absent. The coach must go and find him, and must then await his pleasure. All is in a bad state: the axles are not greased and they grate; the wheels are badly fixed; the shaft dangles dangerously; the horses, although of noble race, are dirty and ill-fed; the harness is worn and the reins are not strong. The coachman is asleep: his hands have slipped to his knees and hardly hold the reins, which can fall from them at any moment.

The coach nevertheless continues to move forward, but does so in a way which presages no happiness. Abandoning the road, it is rolling down the slope in such a way that the coach is now pushing the horses, which are unable to hold it back. The coachman, fallen into a deep sleep, is swaying in his seat at risk of falling off. Obviously a sad fate awaits such a coach.

This image provides a highly appropriate analogy for the condition of most men, and it is worth taking as an object of meditation. Salvation may however present itself. Another coachman, this one quite awake, may pass by the same route and observe the coach in its sad situation. If he is not much in a hurry, he may perhaps stop to help the coach that is in distress. He will first help the horses hold back the coach from slipping down the slope. Then he will awaken the sleeping driver and together with him will try to bring the coach back to the road. He will lend fodder and money. He might also give advice on the care of the horses, the address of an inn and a coach repairer, and indicate the proper route to follow.

It will be up to the assisted coachman afterward to profit, by his own efforts, from the help and the information received. It will be incumbent on him from this point on to put all things in order and, open eyed, to follow the path he had abandoned.

He will above all fight against sleep, for if he falls asleep again, and if the coach leaves the road again and again finds itself in the same danger, he cannot hope that chance will smile upon him a second time; that another coachman will pass at that moment and at that place and come to his aid once again.

\flright{\textsc{Boris Mouravieff}, \emph{Gnosis}, Book One}



\flrightit{Posted on 2010-04-06 by Cologero}

\begin{center}* * *\end{center}

\begin{footnotesize}
\begin{sffamily}
\texttt{Barb on 2010-04-08 at 04:08 said:}

The meek will inherit the earth... but I wonder what happens to those who let the intellect ``rule''... without hearing the notes of love... or worse yet... misinterpreting those notes -- not hearing the correct melody? Christ talks about those who are meek, -but only those who first seek God -- and walk in a relationship with Him will really ever discover and inherit the kingdom of God.

\hfill

\texttt{Cologero on 2010-04-09 at 09:32 said:}

You are confusing the intellect in this sense --- the ancient and traditional sense --- with the ordinary mind. The twin horses of desire and passion are unruly. unless put under the control of the intellect. Of course, the intellect is oriented toward the whole, toward God, though it is barely aware, typically. Love is desire and passion made self-conscious.

\hfill

\texttt{Roger Buck on 2010-04-17 at 06:46 said:}

I hesitate to comment. Poor web access among other probems makes my assimilation of your very interesting site poor for the moment.

But if you are interested in any initial impressions under the heading of HALF FORMED

I will offer this.

Reading you leaves me strongly with the Second Arcanum of the Seated Woman. The correction the author of a meditation on that Arcanum gives regarding St Yves Alvedrye ( in the DeMaistre stream ... )

A concern the author had in mind I think of prioritising either intellect or power over love, that was so common.

I repeat: I am operating on half formed possibly half baked impressions

Vey, very interesting Charles. You will hear from me privately this weekend ... 

\hfill

\texttt{Cologero on 2010-04-17 at 12:25 said:}

I always appreciate your comments.

I don't believe this post even addressed the question of ``love'', though the topic is certainly planned.

Jesus said, ``if you love me, you will obey my commandments.''\\ There is an awful lot there. Without getting ahead of a planned future post:\\ 1) There must be an intellectual understanding of the good, i.e., the commandments\\ 2) There must be power to obey the commandemts (that is, manifest the Good)

Building on ideas from Solovyov: for the pagan (Plato), it was enough to contemplate the Good, in effect, as object.\\ The Christian idea is that the Good is also Subject. A subject involves the Will, and Love is to Will the Good, not just contemplate it. So, yes, Love is also Power. But it is Power conscious of itself, whereas the Will-to-Power in itself is blind and unconscious.

That is what I am getting at in the meaning of ``meekness'' ... the integral person must not only Know the Good above him, he must have power over what is below him, in order to ``obey'' that Good that he knows. According to Tomberg, matter is objectified Will, so the issue of Power is always present. The question is whose Will? Is it a conscious Willing of the Good, which is Love? Or something spontaneous, deterministic, and blind?

\end{sffamily}
\end{footnotesize}

